%==============================================================================
% Praca dyplomowa (UR) — plik główny
%==============================================================================

\documentclass{urdpl}

%------------------------------------------------------------------------------
% JĘZYK / KODOWANIE (pdfLaTeX)
%------------------------------------------------------------------------------
\usepackage[main=polish,english]{babel}
\usepackage[utf8]{inputenc} % tylko pdfLaTeX

%------------------------------------------------------------------------------
% MATEMATYKA
%------------------------------------------------------------------------------
\usepackage{mathtools}
\usepackage{amssymb}

%------------------------------------------------------------------------------
% HYPERLINKI
%------------------------------------------------------------------------------
\usepackage[hidelinks]{hyperref}

%------------------------------------------------------------------------------
% FLOATY
%------------------------------------------------------------------------------
\usepackage{float}
% Trzyma rysunki w obrębie sekcji (zrzut nie ucieka do innego ekranu).
\usepackage[section]{placeins}

%------------------------------------------------------------------------------
% GRAFIKA WEKTOROWA (diagram Gantta, schematy)
%------------------------------------------------------------------------------
\usepackage{tikz}

%------------------------------------------------------------------------------
% TABELE (wg potrzeby)
%------------------------------------------------------------------------------
\usepackage{booktabs}
\usepackage{multirow}
\usepackage{tabularx}
\usepackage{array}
\usepackage{makecell}
\usepackage[flushleft]{threeparttable}

\usepackage{listings}
\usepackage{xcolor}

\newcolumntype{C}[1]{>{\hsize=#1\hsize\centering\arraybackslash}X}

%------------------------------------------------------------------------------
% LISTINGI
%------------------------------------------------------------------------------
\usepackage{listings}
\usepackage{xcolor}

\renewcommand{\lstlistlistingname}{Spis listingów}
\renewcommand{\lstlistingname}{Listing}
\providecommand{\listingname}{\lstlistingname} % alias do odwołań w tekście

\definecolor{codegreen}{rgb}{0,0.6,0}
\definecolor{codegray}{rgb}{0.5,0.5,0.5}

\lstdefinestyle{mystyle}{
  basicstyle=\ttfamily\footnotesize,
  breaklines=true,
  numbers=left,
  numberstyle=\tiny\color{codegray},
  numbersep=6pt,
  showstringspaces=false,
  tabsize=2,
  keywordstyle=\color{blue},
  commentstyle=\color{codegreen},
  stringstyle=\color{red},
  literate={ą}{{\k{a}}}1 {ć}{{\'c}}1 {ę}{{\k{e}}}1 {ó}{{\'o}}1 {ń}{{\'n}}1 {ł}{{\l{}}}1 {ś}{{\'s}}1 {ź}{{\'z}}1 {ż}{{\.z}}1
           {Ą}{{\k{A}}}1 {Ć}{{\'C}}1 {Ę}{{\k{E}}}1 {Ó}{{\'O}}1 {Ń}{{\'N}}1 {Ł}{{\L{}}}1 {Ś}{{\'S}}1 {Ź}{{\'Z}}1 {Ż}{{\.Z}}1
}
\lstset{style=mystyle}



\lstdefinestyle{matlabStyle}{
    language=Matlab,
    basicstyle=\ttfamily\small,
    keywordstyle=\color{blue}\bfseries,
    commentstyle=\color{green!50!black},
    stringstyle=\color{red},
    numbers=left,
    numberstyle=\tiny\color{gray},
    stepnumber=1,
    numbersep=10pt,
    % backgroundcolor=\color{gray!10},
    % frame=single,
    breaklines=true,
    showstringspaces=false,
    tabsize=4
}
%------------------------------------------------------------------------------
% PODKREŚLANIE + NUMERY LINII (opcjonalnie)
%------------------------------------------------------------------------------
\usepackage[normalem]{ulem}
\usepackage{lineno}

%------------------------------------------------------------------------------
% CYTOWANIA
%------------------------------------------------------------------------------
\usepackage{csquotes}

%------------------------------------------------------------------------------
% Nazwy podpisów (opcjonalnie; zwykle babel i klasa wystarczą)
%------------------------------------------------------------------------------
\AtBeginDocument{
  \renewcommand{\tablename}{Tabela}
  \renewcommand{\figurename}{Rys.}
}

%------------------------------------------------------------------------------
% ZRZUTY EKRANU / ILUSTRACJE
% Pokazuje obraz, jeśli plik istnieje w figures/, w przeciwnym razie ramkę-
% zaślepkę. Dzięki temu dokument zawsze się kompiluje, a ilustracje pojawiają
% się automatycznie po wrzuceniu plików PNG.
% Użycie: \figscreen{nazwa.png}{szerokosc 0-1}{Podpis.}{etykieta}
%------------------------------------------------------------------------------
\newcommand{\figscreen}[4]{%
  \begin{figure}[htbp]
    \centering
    \IfFileExists{figures/#1}{%
      \includegraphics[width=#2\linewidth,height=0.82\textheight,keepaspectratio]{figures/#1}%
    }{%
      \fbox{\parbox[c][5cm][c]{0.85\linewidth}{\centering\itshape
        Miejsce na ilustrację --- zapisz plik \texttt{\detokenize{figures/#1}}}}%
    }
    \caption{#3\label{#4}}
  \end{figure}%
}

%------------------------------------------------------------------------------
% METADANE PRACY (dla urdpl.cls)
%------------------------------------------------------------------------------
\author{Michał Nocek}
\shortauthor{M. Nocek}
\noAlbum{134951}

\titlePL{TestPortal - aplikacja do zarządzania testami i ocenami}
\titleEN{TestPortal - a test and grade management application}

\shorttitlePL{TestPortal}
\shorttitleEN{TestPortal}

\thesistype{Praca projektowa}
\thesisDone{Praca wykonana pod kierunkiem\\ }
\supervisor{mgr inż. Ewa Żesławska}

\degreeprogramme{Informatyka}

\date{2026}

\department{Instytut Informatyki}
\faculty{Wydział Nauk Ścisłych i Technicznych}


%------------------------------------------------------------------------------
% NUMERACJA
%------------------------------------------------------------------------------
\setlength{\cftsecnumwidth}{10mm}
\setcounter{secnumdepth}{4}
\setcounter{tocdepth}{2}
\brokenpenalty=10000\relax

%==============================================================================
\begin{document}
%==============================================================================


%------------------------------------------------------------------------------
% Strona tytułowa + podziękowania
%------------------------------------------------------------------------------
\titlepages

% (opcjonalnie) styl plain bez numerów na stronach typu plain
\fancypagestyle{plain}{
  \fancyhf{}
  \renewcommand{\headrulewidth}{0pt}
  \renewcommand{\footrulewidth}{0pt}
}

%------------------------------------------------------------------------------
% Streszczenie jest pierwszym, NUMEROWANYM rozdziałem (jak w pracy projektowej)
% — patrz chapters/streszczenie.tex, włączony niżej razem z resztą rozdziałów.
%------------------------------------------------------------------------------

%------------------------------------------------------------------------------
% Spis treści
%------------------------------------------------------------------------------

\cleardoublepage
\selectlanguage{polish}
\tableofcontents

\clearpage
% \linenumbers % jeśli nie chcesz numerów linii w całości, usuń lub przenieś

%------------------------------------------------------------------------------
% Rozdziały
%------------------------------------------------------------------------------
\chapter{Streszczenie w języku polskim i angielskim}
\label{cha:streszczenie}

\section{Streszczenie w języku polskim}
\label{sec:streszczeniePL}

Praca opisuje aplikację \emph{Testownik} - program do tworzenia i przeprowadzania testów oraz
zarządzania wynikami uczniów. Celem projektu było danie nauczycielowi narzędzia, które obejmuje cały
cykl pracy z testem: od ułożenia pytań, przez przypisanie testu grupom uczniów, po automatyczne
i ręczne ocenianie odpowiedzi oraz analizę wyników.

Aplikacja została napisana w języku C\# na platformie .NET z wykorzystaniem technologii
\emph{Windows Presentation Foundation} (WPF) do budowy interfejsu graficznego oraz \emph{Entity
Framework Core} jako warstwy dostępu do relacyjnej bazy danych PostgreSQL. System rozróżnia dwie
role użytkowników - nauczyciela i ucznia - udostępniając każdej z nich odrębny zestaw funkcji.
Nauczyciel tworzy testy złożone z pytań jednokrotnego wyboru, wielokrotnego wyboru oraz pytań
otwartych, zarządza grupami uczniów i analizuje wyniki, natomiast uczeń rozwiązuje przypisane mu
testy w ograniczonym czasie i przegląda swoje oceny. W aplikacji zaimplementowano między innymi
limit podejść do testu, automatyczne ocenianie pytań zamkniętych, ręczne ocenianie i edycję oceny,
sterowanie widocznością testów, wykrywanie prób oszustwa, eksport wyników do pliku CSV oraz
haszowanie haseł. Efektem końcowym jest działająca aplikacja okienkowa gotowa do użycia.

\section{Streszczenie w języku angielskim}
\label{sec:streszczenieEN}

\selectlanguage{english}
This document describes the \emph{Testownik} application - an information system for creating,
sharing and conducting tests, as well as for managing students' results. The main goal of the
project was to provide the teacher with a tool that automates the entire workflow of working with a
test: from composing questions, through assigning the test to groups of students, up to automatic
and manual grading of answers and the statistical analysis of results.

The application was written in C\# on the .NET platform using \emph{Windows Presentation
Foundation} (WPF) for the graphical user interface and \emph{Entity Framework Core} as the data
access layer for a relational PostgreSQL database. The system distinguishes two user roles - the
teacher and the student - providing each of them with a separate set of features. The teacher
creates tests consisting of single-choice, multiple-choice and open questions, manages groups of
students and analyses results, while the student solves the assigned tests within a limited time
and reviews their grades. The application implements, among others, a limit on the number of
attempts, automatic grading of closed questions, manual grading and grade editing, control over test
visibility, cheating detection, exporting results to a CSV file and password hashing. The end result
is a working desktop application ready for use.
\selectlanguage{polish}
     % 1
\chapter{Wprowadzenie i analiza problemu}
\label{cha:wprowadzenie}

\section{Opis analizowanego problemu}
\label{sec:problem}

Sprawdzanie wiedzy uczniów za pomocą testów jest jednym z podstawowych zadań nauczyciela, jednak
w tradycyjnej, papierowej formie jest ono czasochłonne i mało wygodne. Wymaga ręcznego przygotowania
arkuszy, ich powielenia, a następnie żmudnego sprawdzania odpowiedzi i podliczania punktów. Dodatkowo
gromadzenie wyników oraz analiza tego, które zagadnienia sprawiają uczniom najwięcej trudności, są
w takim modelu praktycznie niemożliwe do przeprowadzenia na bieżąco.

Istnieją platformy do testów online, jednak najczęściej są to rozwiązania webowe, wymagające stałego
dostępu do Internetu i kont w usługach zewnętrznych, a ich pełne możliwości są zwykle płatne. Brakuje
prostego, samodzielnego narzędzia desktopowego, które łączyłoby tworzenie testów, ich automatyczne
ocenianie oraz analizę wyników w jednym miejscu, z danymi przechowywanymi w lokalnej bazie.

\section{Cel projektu}
\label{sec:cel}

Celem projektu było stworzenie od podstaw funkcjonalnej aplikacji okienkowej w języku C\#,
współpracującej z relacyjną bazą danych, która realizuje pełny cykl pracy z testem:

\begin{itemize}
  \item rejestrację i logowanie użytkowników z bezpiecznym przechowywaniem haseł oraz rozróżnieniem
        ról nauczyciela i ucznia,
  \item tworzenie testów złożonych z pytań jednokrotnego wyboru, wielokrotnego wyboru i pytań
        otwartych, z możliwością sterowania ich widocznością,
  \item przypisywanie testów grupom uczniów,
  \item rozwiązywanie testów przez uczniów w ograniczonym czasie, z obsługą wielu podejść,
  \item automatyczne ocenianie pytań zamkniętych oraz ręczne ocenianie i korektę punktów,
  \item wykrywanie prób oszustwa podczas rozwiązywania testu,
  \item analizę wyników w postaci statystyk oraz eksport do pliku CSV.
\end{itemize}

\section{Zakres realizowanych funkcjonalności}
\label{sec:zakres}

Zakres projektu obejmuje następujące moduły:

\begin{enumerate}
  \item \textbf{Moduł uwierzytelniania} - rejestracja (z wyborem roli), logowanie, haszowanie haseł,
        rozróżnienie ról nauczyciel/uczeń.
  \item \textbf{Moduł zarządzania grupami} - tworzenie grup, zarządzanie ich składem, przynależność
        ucznia do wielu grup.
  \item \textbf{Moduł testów} - tworzenie, edycja, duplikowanie i usuwanie testów, edytor pytań,
        ustawienia (czas, próg, liczba podejść), sterowanie widocznością (aktywny/nieaktywny).
  \item \textbf{Moduł rozwiązywania testów} - rozwiązywanie w ograniczonym czasie, nawigacja między
        pytaniami, limit podejść, wykrywanie opuszczenia okna testu.
  \item \textbf{Moduł oceniania} - automatyczne ocenianie pytań zamkniętych, ręczne ocenianie pytań
        otwartych oraz korekta oceny dowolnego podejścia.
  \item \textbf{Moduł statystyk} - zdawalność, średni wynik, trudność pytań, tabela wyników,
        eksport do CSV.
\end{enumerate}

\section{Przegląd istniejących rozwiązań}
\label{sec:przeglad}

Przed realizacją projektu przeanalizowano kilka istniejących narzędzi o podobnym przeznaczeniu.

\subsection{Kahoot!}

Popularna platforma do quizów i testów, działająca w przeglądarce, często wykorzystywana w szkołach.
Jej zaletą jest grywalizacja i prostota, jednak jest to wyłącznie usługa internetowa, wymagająca
stałego dostępu do sieci i konta, a obsługuje głównie pytania zamknięte. Pełne funkcje (m.in.
rozbudowane raporty) są dostępne w wersji płatnej.

\subsection{Microsoft Forms / Google Forms}

Webowe narzędzia do tworzenia formularzy i quizów, umożliwiające automatyczne sprawdzanie odpowiedzi
zamkniętych i zbieranie wyników. Są wygodne i darmowe w podstawowym zakresie, ale stanowią ogólne
narzędzia do ankiet - brakuje im funkcji typowo dydaktycznych, takich jak limit podejść, grupy
uczniów z przypisanymi testami, ręczne ocenianie pytań otwartych z korektą czy wykrywanie prób
oszustwa. Wymagają też konta w usłudze chmurowej i połączenia z Internetem.

\subsection{Moduł testów w systemie Moodle}

Moodle to rozbudowana platforma e-learningowa oferująca zaawansowany moduł testów (różne typy pytań,
banki pytań, statystyki). Jest jednak ciężka we wdrożeniu - wymaga serwera, bazy danych i konfiguracji
całej platformy - co czyni ją nieproporcjonalnie złożoną do prostego scenariusza, w którym nauczyciel
chce jedynie szybko ułożyć test i sprawdzić wyniki na własnym komputerze.

\subsection{Wnioski z analizy}

Żadne z analizowanych rozwiązań nie łączy lekkiej, samodzielnej aplikacji desktopowej z własną,
lokalną bazą danych, automatycznym i ręcznym ocenianiem (w tym pytań otwartych), obsługą grup,
limitem podejść, sterowaniem widocznością testu oraz wykrywaniem prób oszustwa - bez konieczności
korzystania z usług chmurowych. Te cechy wyznaczyły zakres tego projektu.
     % 2
\chapter{Wymagania systemowe}
\label{cha:wymagania}

\section{Wymagania funkcjonalne}
\label{sec:wymaganiaFunkcjonalne}

\begin{itemize}
  \item \textbf{Zarządzanie kontami i rolami użytkowników}
  \begin{itemize}
    \item Użytkownik może samodzielnie zarejestrować nowe konto (imię, nazwisko, login, hasło),
          wybierając przy rejestracji rolę - ucznia lub nauczyciela. Login jest unikalny w skali
          całego systemu.
    \item System rozróżnia dwie role - nauczyciela i ucznia - i po zalogowaniu udostępnia każdej
          z nich odrębny zestaw funkcji oraz inny układ nawigacji. Każdy nauczyciel zarządza
          wyłącznie własnymi testami.
  \end{itemize}

  \item \textbf{Zarządzanie grupami uczniów}
  \begin{itemize}
    \item Nauczyciel może tworzyć nowe grupy opatrzone nazwą i opisem oraz zarządzać ich składem.
    \item Jeden uczeń może należeć do wielu grup jednocześnie.
  \end{itemize}

  \item \textbf{Tworzenie i edycja testów}
  \begin{itemize}
    \item Nauczyciel tworzy testy złożone z pytań trzech typów: jednokrotnego wyboru, wielokrotnego
          wyboru oraz pytań otwartych.
    \item Dla każdego testu określa nazwę, przedmiot, limit czasu, próg zaliczenia, dozwoloną liczbę
          podejść, widoczność oraz grupy, którym test jest udostępniany.
    \item Test można edytować, duplikować i usuwać. Po pojawieniu się rozwiązań edycja pytań jest
          blokowana, aby nie naruszyć powiązanych z nimi odpowiedzi.
    \item Tylko aktywny test jest widoczny dla uczniów; nowy test powstaje jako nieaktywny (szkic).
  \end{itemize}

  \item \textbf{Rozwiązywanie testów}
  \begin{itemize}
    \item Uczeń widzi listę testów przypisanych do jego grup, z podziałem na rozwiązane
          i nierozwiązane.
    \item Podczas rozwiązywania działa licznik czasu; po jego upływie test jest automatycznie
          oddawany. Liczba podejść jest ograniczona wartością ustawioną przez nauczyciela.
  \end{itemize}

  \item \textbf{Ocenianie odpowiedzi}
  \begin{itemize}
    \item Pytania zamknięte są oceniane automatycznie, w wielokrotnym wyborze z punktami częściowymi.
    \item Pytania otwarte ocenia ręcznie nauczyciel; może też skorygować punkty dowolnego pytania
          każdego podejścia. Przy wielu podejściach liczy się najlepszy wynik.
  \end{itemize}

  \item \textbf{Wykrywanie prób oszustwa}
  \begin{itemize}
    \item Opuszczenie okna testu jest wykrywane: pierwsze skutkuje ostrzeżeniem, kolejne -
          zakończeniem testu z wynikiem 0\% i oznaczeniem próby oszustwa.
  \end{itemize}

  \item \textbf{Statystyki i eksport wyników}
  \begin{itemize}
    \item Nauczyciel ma dostęp do statystyk testu (zdawalność, średni wynik, najtrudniejsze pytanie,
          poprawność pytań) oraz może wyeksportować wyniki do pliku CSV.
  \end{itemize}
\end{itemize}

\section{Wymagania niefunkcjonalne}
\label{sec:wymaganiaNiefunkcjonalne}

\textbf{Bezpieczeństwo}
\begin{itemize}
  \item Dostęp do funkcji systemu jest możliwy wyłącznie po zalogowaniu, a zakres uprawnień zależy
        od roli. Hasła nie są przechowywane jawnie - zapisuje się ich skrót (SHA-256). Unikalność
        loginów jest wymuszana na poziomie bazy danych.
\end{itemize}

\textbf{Niezawodność i obsługa błędów}
\begin{itemize}
  \item Operacje na bazie danych oraz operacje wejścia/wyjścia są zabezpieczone wspólnym mechanizmem
        obsługi wyjątków. W razie błędu użytkownik otrzymuje zrozumiały komunikat, a aplikacja nie
        ulega awarii.
\end{itemize}

\textbf{Spójność danych}
\begin{itemize}
  \item Integralność jest wymuszana kluczami obcymi z odpowiednimi regułami usuwania, unikalnym
        indeksem oraz mechanizmem migracji wersjonującym schemat bazy.
\end{itemize}

\textbf{Użyteczność}
\begin{itemize}
  \item Interfejs jest czytelny i spójny, nawigacja odbywa się przez stały panel boczny, a komunikaty
        prezentowane są w języku polskim. Nowy użytkownik bez trudu odnajduje podstawowe funkcje.
\end{itemize}

\textbf{Wydajność}
\begin{itemize}
  \item Kluczowe operacje (logowanie, wczytywanie testów, zapis wyniku) wykonują się w czasie
        nieprzekraczającym kilku sekund.
\end{itemize}

\section{Założenia dotyczące działania systemu}
\label{sec:zalozeniaSystemu}

Przyjęto następujące założenia:

\begin{itemize}
  \item Aplikacja działa w systemie Windows 10/11 z zainstalowanym środowiskiem uruchomieniowym .NET.
  \item Dostępna jest instancja serwera PostgreSQL; baza danych i jej schemat są tworzone automatycznie
        przy pierwszym uruchomieniu za pomocą migracji Entity Framework Core.
  \item Dane początkowe (przykładowy nauczyciel, uczniowie, grupy i testy) są wstawiane automatycznie,
        co umożliwia natychmiastowe przetestowanie aplikacji.
  \item Każdy nauczyciel pracuje na własnym zbiorze testów; grupy uczniów są zasobem współdzielonym.
\end{itemize}
        % 3
\chapter{Architektura i struktura aplikacji}
\label{cha:architektura}

\section{Architektura aplikacji}
\label{sec:architektura}

Aplikacja ma architekturę warstwową. Wyróżniono cztery warstwy o odrębnych
odpowiedzialnościach:

\begin{enumerate}
  \item \textbf{Warstwa modelu (\texttt{Models})} - klasy encji odwzorowujące obiekty dziedziny
        (użytkownik, grupa, test, pytanie, podejście, odpowiedź), mapowane na tabele bazy danych
        przez Entity Framework Core.
  \item \textbf{Warstwa dostępu do danych (\texttt{Data})} - kontekst bazy danych
        (\texttt{AppDbContext}) z konfiguracją relacji i dziedziczenia, mechanizm migracji oraz klasa
        zasilająca bazę danymi początkowymi (\texttt{DbSeeder}).
  \item \textbf{Warstwa logiki i usług (\texttt{Services})} - repozytorium operacji bazodanowych
        (\texttt{DbRepository}), serwis oceniania (\texttt{GradingService}), eksport do CSV
        (\texttt{CsvExporter}), haszowanie haseł (\texttt{PasswordHash}) oraz wspólna obsługa
        wyjątków (\texttt{Safe}). Tu znajdują się także obiekty transferu danych (DTO).
  \item \textbf{Warstwa prezentacji (\texttt{Views})} - okna i ekrany aplikacji zbudowane w
        technologii WPF, odpowiedzialne za wyświetlanie danych i obsługę interakcji użytkownika.
\end{enumerate}

Przepływ jest jednokierunkowy: widok wywołuje warstwę usług, ta korzysta z repozytorium, a
repozytorium komunikuje się z bazą przez kontekst Entity Framework Core. Widoki nigdy nie odwołują
się do bazy bezpośrednio - zawsze pośredniczy w tym warstwa usług.

\section{Struktura folderów i organizacja kodu}
\label{sec:foldery}

Kod źródłowy podzielono na foldery odpowiadające warstwom aplikacji (listing~\ref{lst:foldery}).

\begin{lstlisting}[caption={Struktura katalogów projektu.}, label=lst:foldery]
ProjektPO2/
|-- App.xaml(.cs)              // punkt startowy aplikacji
|-- MainWindow.xaml(.cs)       // powloka: panel boczny + obszar tresci
|-- Theme.xaml                 // motyw: kolory, style kontrolek
|-- Models/                    // encje (User, Quiz, Question, ...)
|-- Data/                      // AppDbContext, DbSeeder, fabryka kontekstu
|-- Migrations/                // migracje EF Core
|-- Services/                  // DbRepository, GradingService, CsvExporter,
|                              //   PasswordHash, Safe, DTO (AppModels)
|-- Views/                     // ekrany aplikacji (Login, Teacher*, Quiz, ...)
\-- Commands/                  // RelayCommand
\end{lstlisting}

\section{Zastosowane wzorce projektowe}
\label{sec:wzorce}

W projekcie wykorzystano kilka rozwiązań i wzorców porządkujących kod:

\begin{itemize}
  \item \textbf{Architektura warstwowa} - oddzielenie modelu, dostępu do danych, logiki i prezentacji.
  \item \textbf{Repozytorium (Repository)} - klasa \texttt{DbRepository} hermetyzuje cały dostęp do
        bazy danych i odcina od niej resztę aplikacji.
  \item \textbf{Obiekty transferu danych (DTO)} - lekkie klasy w warstwie usług oddzielają dane
        prezentowane w widokach od encji bazodanowych.
  \item \textbf{Dziedziczenie i polimorfizm} - trzy hierarchie klas abstrakcyjnych (użytkownicy,
        pytania, odpowiedzi) mapowane na bazę strategią Table-per-Hierarchy.
  \item \textbf{Interfejs \texttt{IPublishable}} - implementowany przez klasę \texttt{Quiz},
        ujednolica sterowanie widocznością testu (metody \texttt{Publish()} i \texttt{Unpublish()}).
\end{itemize}

\section{Mechanizmy walidacji danych}
\label{sec:walidacja}

Walidacja realizowana jest na trzech poziomach:

\begin{enumerate}
  \item \textbf{Interfejs użytkownika} - pola formularzy są sprawdzane przed zapisem: niepuste
        wymagane pola, liczby ograniczane do dozwolonego zakresu (np. czas, próg, liczba podejść),
        a edytor testu blokuje zapis testu niekompletnego (pytanie bez poprawnej odpowiedzi, puste
        pola, brak przypisanej grupy).
  \item \textbf{Logika aplikacji} - serwis oceniania pilnuje zakresu przyznawanych punktów, a
        repozytorium dba o spójność operacji (np. bezpieczną kolejność usuwania powiązanych danych).
  \item \textbf{Baza danych} - ograniczenia na poziomie tabel (klucze obce, unikalny indeks)
        zapewniają spójność nawet przy bezpośrednim dostępie do bazy.
\end{enumerate}

\section{Obsługa błędów i wyjątków}
\label{sec:bledy}

Operacje na bazie danych oraz operacje wejścia/wyjścia są wykonywane za pośrednictwem wspólnej klasy
\texttt{Safe}, która przechwytuje wyjątki i prezentuje użytkownikowi czytelny komunikat zamiast
powodować awarię aplikacji. Dzięki temu np. brak połączenia z bazą czy błąd zapisu pliku CSV nie
zatrzymuje programu, a użytkownik dowiaduje się o przyczynie problemu.
     % 4
\chapter{Technologie wykorzystane w projekcie}
\label{cha:technologie}

\section{Technologie interfejsu użytkownika}
\label{sec:techUI}

Warstwę prezentacji zrealizowano w technologii \emph{Windows Presentation Foundation}
(WPF)~\cite{msdocsWPF} na platformie .NET, w języku C\#~\cite{msdocsCSharp}. WPF posłużył do budowy
okien aplikacji, paska bocznego nawigacji oraz ekranów poszczególnych modułów. Wygląd aplikacji
(kolory, style kontrolek) zdefiniowano w słowniku zasobów \texttt{Theme.xaml}, co oddziela warstwę
wizualną od logiki. Strona logowania zbudowana jest deklaratywnie w pliku XAML, natomiast pozostałe
ekrany tworzone są dynamicznie w kodzie, z wykorzystaniem wspólnych metod pomocniczych budujących
powtarzalne elementy interfejsu (karty, znaczniki statusu, wykresy).

\section{Technologie dostępu do danych}
\label{sec:techDane}

Warstwę dostępu do danych oparto na bibliotece \emph{Entity Framework Core}~\cite{msdocsEFCore} -
mechanizmie mapowania obiektowo-relacyjnego (ORM). Pozwala on operować na bazie za pomocą obiektów
języka C\# i zapytań LINQ, bez ręcznego pisania większości zapytań SQL. Systemem zarządzania bazą danych jest
\emph{PostgreSQL}~\cite{postgresql}, łączony z aplikacją za pośrednictwem dostawcy
\emph{Npgsql}~\cite{npgsql}. Schemat bazy jest wersjonowany i tworzony automatycznie poprzez migracje
Entity Framework Core.

\section{Dodatkowe biblioteki i narzędzia}
\label{sec:techDodatkowe}

W projekcie wykorzystano ponadto:

\begin{itemize}
  \item \texttt{System.Security.Cryptography} - moduł platformy .NET użyty do wyznaczania skrótu
        SHA-256 haseł użytkowników.
  \item \texttt{System.IO} (\texttt{StreamWriter}) - strumieniowy zapis pliku CSV z wynikami testu,
        z kodowaniem UTF-8 obsługującym polskie znaki.
  \item \emph{Visual Studio} - środowisko programistyczne wykorzystane do budowy i uruchamiania
        aplikacji.
  \item \emph{pgAdmin} - narzędzie do zarządzania bazą danych PostgreSQL.
  \item \emph{Git} - system kontroli wersji projektu.
\end{itemize}

\section{Uzasadnienie wyboru technologii}
\label{sec:techUzasadnienie}

WPF, Entity Framework Core i PostgreSQL dobrano pod charakter aplikacji:

\begin{itemize}
  \item \textbf{WPF} wybrano, bo aplikacja jest programem desktopowym z rozbudowanym interfejsem
        (karty, wykresy, stylizacja), a nie potrzebuje serwera ani przeglądarki.
  \item \textbf{Entity Framework Core} wybrano zamiast bezpośredniego ADO.NET, ponieważ model danych
        opiera się na dziedziczeniu (trzy hierarchie klas) i wielu relacjach. ORM odwzorowuje je
        automatycznie (strategią Table-per-Hierarchy), a migracje pozwalają wersjonować schemat bez
        ręcznego pisania skryptów SQL, co skraca kod warstwy danych.
  \item \textbf{PostgreSQL} to darmowy, otwarty i niezawodny system bazodanowy, dobrze wspierany
        przez Entity Framework Core poprzez dostawcę Npgsql.
\end{itemize}
      % 5
\chapter{Projekt bazy danych}
\label{cha:baza}

\section{Opis modelu danych}
\label{sec:modelDanych}

Wszystkie informacje o użytkownikach, grupach, testach, pytaniach, podejściach i odpowiedziach są
przechowywane w relacyjnej bazie danych PostgreSQL. Najważniejsze tabele przedstawiono
w tabeli~\ref{tab:tabele}. Można je podzielić na trzy obszary: konta użytkowników i grupy, testy
wraz z pytaniami oraz podejścia uczniów wraz z ich odpowiedziami.

\begin{table}[htbp]
  \centering
  \caption{Najważniejsze tabele bazy danych aplikacji.}
  \label{tab:tabele}
  \begin{tabularx}{\linewidth}{>{\bfseries}l X}
    \toprule
    Tabela & Przeznaczenie \\
    \midrule
    Users & Konta użytkowników (wspólna tabela nauczycieli i uczniów, rozróżnianych
            dyskryminatorem); login, skrót hasła, imię i nazwisko. \\
    Groups & Grupy uczniów (nazwa, opis). \\
    StudentGroups & Tabela łącząca uczniów z grupami (relacja wiele-do-wielu). \\
    Quizzes & Testy (nazwa, przedmiot, limit czasu, próg, liczba podejść, widoczność). \\
    Questions & Pytania (wspólna tabela pytań zamkniętych i otwartych); treść, punkty, kolejność. \\
    AnswerOptions & Warianty odpowiedzi pytań zamkniętych; treść, flaga poprawności. \\
    TestGroups & Tabela łącząca testy z grupami (relacja wiele-do-wielu). \\
    QuizAttempts & Podejścia uczniów do testów; wynik, procent, status, czas, data. \\
    StudentAnswers & Odpowiedzi ucznia (wspólna tabela odpowiedzi zamkniętych i otwartych). \\
    SelectedOptions & Warianty zaznaczone przez ucznia w odpowiedzi zamkniętej. \\
    \bottomrule
  \end{tabularx}
\end{table}

\section{Diagram ERD}
\label{sec:erd}

Centralnymi tabelami modelu są \texttt{Quizzes} (testy) oraz \texttt{Users} (użytkownicy). Test
posiada pytania i podejścia, a podejście ucznia gromadzi jego odpowiedzi. Pełny schemat powiązań
przedstawia diagram ERD na rysunku~\ref{rys:erd}.

\figscreen{erd.png}{0.95}{Diagram ERD bazy danych aplikacji.}{rys:erd}

\section{Opis relacji i kluczy}
\label{sec:relacje}

Każda tabela posiada klucz główny typu całkowitego (\texttt{IDENTITY}). Główne relacje:

\begin{itemize}
  \item \texttt{Quizzes} - \texttt{Questions} oraz \texttt{Quizzes} - \texttt{QuizAttempts}: relacje
        jeden-do-wielu (test ma wiele pytań i wiele podejść).
  \item \texttt{ClosedQuestion} - \texttt{AnswerOptions}: pytanie zamknięte ma wiele wariantów
        odpowiedzi.
  \item \texttt{QuizAttempts} - \texttt{StudentAnswers} - \texttt{SelectedOptions}: podejście ma wiele
        odpowiedzi, a odpowiedź zamknięta - wiele zaznaczonych wariantów.
  \item \texttt{StudentGroups} łączy uczniów i grupy w relacji wiele-do-wielu (klucz złożony
        \texttt{StudentId} + \texttt{GroupId}); analogicznie \texttt{TestGroups} łączy testy z grupami.
\end{itemize}

\section{Normalizacja}
\label{sec:normalizacja}

Baza danych spełnia trzecią postać normalną (3NF) - każda tabela przechowuje dane dotyczące jednego
pojęcia, wartości w kolumnach są atomowe, a powtarzające się informacje zostały zastąpione kluczami
obcymi. Relacje wiele-do-wielu (uczeń-grupa oraz test-grupa) rozwiązano za pomocą tabel łączących
(\texttt{StudentGroups}, \texttt{TestGroups}), co eliminuje powielanie danych. Dziedziczenie klas
odwzorowano strategią Table-per-Hierarchy, w której cała hierarchia mieści się w jednej tabeli,
a typ rekordu wskazuje kolumna dyskryminatora.

\section{Migracje bazy danych}
\label{sec:migracje}

Zmiany schematu bazy są zapisane w postaci migracji Entity Framework Core w katalogu
\texttt{Migrations}. Przy starcie aplikacja wywołuje metodę \texttt{Database.Migrate()}, która tworzy
bazę i nakłada wszystkie migracje. Kolejne migracje wprowadzały m.in. obsługę liczby podejść, zmianę
typu daty oddania na pełny znacznik czasu, unikalny indeks loginu oraz haszowanie haseł istniejących
kont. Dzięki migracjom schemat jest wersjonowany, a baza odtwarzalna bez ręcznego pisania skryptów SQL.

\section{Mechanizmy zapewniające integralność danych}
\label{sec:integralnosc}

\begin{itemize}
  \item \textbf{Integralność referencyjna} - klucze obce na wszystkich powiązaniach. Tam, gdzie dane
        podrzędne nie mają sensu bez nadrzędnych, zastosowano usuwanie kaskadowe (usunięcie testu
        usuwa jego pytania i podejścia); tam, gdzie usunięcie mogłoby naruszyć już oddane odpowiedzi,
        zastosowano regułę ograniczającą, a operacje usuwania wykonywane są we właściwej kolejności
        (najpierw rekordy zależne).
  \item \textbf{Unikalność} - unikalny indeks na loginie użytkownika gwarantuje brak dwóch kont o tej
        samej nazwie.
  \item \textbf{Spójność dziedzinowa} - walidacja po stronie aplikacji nie dopuszcza do zapisu testu
        niekompletnego (np. pytania bez poprawnej odpowiedzi).
\end{itemize}
             % 6
\chapter{Implementacja aplikacji}
\label{cha:implementacja}

\section{Najważniejsze moduły systemu}
\label{sec:moduly}

Logika aplikacji skupiona jest w warstwie usług. Kluczowe klasy to:

\begin{itemize}
  \item \texttt{DbRepository} - centralne repozytorium operacji na bazie danych; pobieranie
        i zapisywanie testów, grup, podejść i wyników oraz mapowanie encji na obiekty DTO,
  \item \texttt{GradingService} - logika oceniania (punkty za pojedyncze pytanie oraz przeliczanie
        całego wyniku podejścia),
  \item \texttt{CsvExporter} - eksport wyników testu do pliku CSV,
  \item \texttt{PasswordHash} - haszowanie haseł,
  \item \texttt{Safe} - wspólna obsługa wyjątków.
\end{itemize}

Wzajemne powiązania najważniejszych klas modelu przedstawia diagram UML na rysunku~\ref{rys:uml}.

\figscreen{uml.png}{0.95}{Diagram UML klas modelu aplikacji.}{rys:uml}

\section{Mechanizm logowania i autoryzacji}
\label{sec:logowanie}

Hasła nie są przechowywane jawnie - zapisywany jest ich skrót SHA-256 w postaci szesnastkowej.
Podczas logowania wyznaczany jest skrót wpisanego hasła i porównywany ze skrótem w bazie
(listing~\ref{lst:hash}). Autoryzacja oparta jest na roli użytkownika: po zalogowaniu obiekt
użytkownika trafia do okna głównego, które buduje inny zestaw modułów dla nauczyciela, a inny dla
ucznia.

\begin{lstlisting}[caption={Haszowanie hasła i logowanie.}, label=lst:hash]
public static string Hash(string password)
{
    byte[] bytes = SHA256.HashData(Encoding.UTF8.GetBytes(password));
    return Convert.ToHexString(bytes).ToLowerInvariant();
}

// logowanie - porownanie skrotu wpisanego hasla ze skrotem w bazie
var hash = PasswordHash.Hash(password);
var user = ctx.Users.FirstOrDefault(
    x => x.Username == username && x.Password == hash);
\end{lstlisting}

\section{Realizacja operacji CRUD}
\label{sec:crud}

Operacje tworzenia, odczytu, modyfikacji i usuwania (CRUD) realizowane są przez repozytorium
\texttt{DbRepository} z użyciem Entity Framework Core. Przykładem nietrywialnej operacji jest
usuwanie testu (listing~\ref{lst:delete}). Ponieważ odpowiedzi uczniów są powiązane z pytaniami
kluczem obcym w trybie ograniczającym (\texttt{Restrict}), bezpośrednie usunięcie testu mogłoby
naruszyć integralność. Dlatego najpierw usuwane są odpowiedzi i podejścia, a dopiero potem sam test -
pozostałe dane (pytania, opcje, przypisania grup) znikają kaskadowo.

\begin{lstlisting}[caption={Bezpieczne usuwanie testu (zachowanie integralności).}, label=lst:delete]
public static void DeleteTest(int id)
{
    using var ctx = Ctx();
    var quiz = ctx.Quizzes
        .Include(q => q.Attempts).ThenInclude(a => a.Answers)
        .FirstOrDefault(q => q.Id == id);
    if (quiz == null) return;

    foreach (var a in quiz.Attempts)
        ctx.StudentAnswers.RemoveRange(a.Answers);
    ctx.QuizAttempts.RemoveRange(quiz.Attempts);
    ctx.Quizzes.Remove(quiz);
    ctx.SaveChanges();
}
\end{lstlisting}

\section{Logika oceniania}
\label{sec:logikaOceniania}

Najważniejszym elementem logiki jest serwis oceniania. Metoda \texttt{Grade} wyznacza liczbę punktów za pojedyncze
pytanie (listing~\ref{lst:grade}). Pierwszeństwo ma ręcznie przyznana ocena - dotyczy to każdego
typu pytania, dzięki czemu nauczyciel może nadpisać również automatyczną ocenę pytania zamkniętego.
W przypadku braku ręcznej oceny pytania zamknięte są liczone automatycznie, przy czym w pytaniach
wielokrotnego wyboru przyznawane są punkty częściowe (proporcjonalnie do liczby trafień
pomniejszonych o błędne wskazania). Pytania otwarte bez ręcznej oceny oczekują na nauczyciela.
Cały wynik podejścia przeliczany jest następnie metodą \texttt{Recompute}, a przy wielu podejściach
do oceny brany jest najlepszy wynik ucznia.

\begin{lstlisting}[caption={Ocenianie pojedynczego pytania (GradingService).}, label=lst:grade]
public static (int pts, bool ok, bool partial, bool pending) Grade(
    Question q, object? ans, IReadOnlyDictionary<int, int>? manual = null)
{
    // reczne nadpisanie punktow ma pierwszenstwo (kazdy typ pytania)
    if (manual != null && manual.TryGetValue(q.Id, out var mo))
        return (mo, mo >= q.Points, mo > 0 && mo < q.Points, false);

    switch (q.Type)
    {
        case QuestionType.Multi:
            var chosen = AsIndices(ans);
            int good = chosen.Count(i => q.Correct.Contains(i));
            int bad  = chosen.Count(i => !q.Correct.Contains(i));
            double raw = q.Correct.Count == 0 ? 0
                         : Math.Max(0, good - bad) / (double)q.Correct.Count;
            int pts = (int)Math.Round(raw * q.Points);
            return (pts, pts == q.Points && q.Points > 0,
                    pts > 0 && pts < q.Points, false);
        // ... pytania jednokrotnego wyboru i otwarte
    }
    return (0, false, false, false);
}
\end{lstlisting}

\section{Organizacja nawigacji}
\label{sec:nawigacja}

Główne okno (\texttt{MainWindow}) pełni rolę powłoki: zawiera stały panel boczny z nawigacją oraz
obszar treści, w którym wyświetlane są kolejne ekrany modułów. Po kliknięciu pozycji menu tworzona
jest odpowiednia strona i wstawiana do obszaru treści. Ekran rozwiązywania testu oraz ekran wyniku
wyświetlane są w trybie pełnoekranowym (bez panelu bocznego), aby skupić uwagę ucznia na zadaniu.

\section{Interfejs użytkownika i stylizacja}
\label{sec:stylizacja}

Wygląd aplikacji zdefiniowano w słowniku zasobów \texttt{Theme.xaml}, obejmującym paletę kolorów oraz
style przycisków, pól i innych kontrolek. Powtarzalne elementy interfejsu (karty, znaczniki statusu,
wykresy kołowe i słupkowe, awatary) budowane są przez wspólne metody pomocnicze, dzięki czemu ekrany
wyglądają jednolicie. Okno główne otwiera się zmaksymalizowane, a układ oparty na siatkach dopasowuje
się do rozmiaru okna.
    % 7
\chapter{Prezentacja warstwy użytkowej projektu}
\label{cha:prezentacja}

W tym rozdziale przedstawiono wizualną stronę aplikacji za pomocą zrzutów ekranu, ilustrujących
kluczowe widoki oraz interakcje użytkownika. Celem jest zaprezentowanie funkcjonalności systemu
z perspektywy użytkownika końcowego - zarówno nauczyciela, jak i ucznia. Kolejne podrozdziały
odpowiadają typowemu scenariuszowi pracy: nauczyciel loguje się, tworzy i udostępnia test, uczeń
rozwiązuje przypisany mu test, a następnie nauczyciel ocenia odpowiedzi i analizuje wyniki.

% =====================================================================
\section{Ekran logowania i rejestracji}
\label{sec:logowanie}

Pierwszym ekranem, z którym styka się użytkownik, jest ekran logowania, umożliwiający dostęp do
systemu zarejestrowanym użytkownikom lub utworzenie nowego konta (ucznia lub nauczyciela). Interfejs
przedstawiono na rysunku~\ref{rys:logowanie}. Składa się on z:

\begin{itemize}
  \item pola \emph{Login}, służącego do wprowadzenia nazwy użytkownika,
  \item pola \emph{Hasło}, zapewniającego maskowanie wpisywanego hasła,
  \item przycisku \emph{Zaloguj się}, który weryfikuje podane dane z bazą danych i w przypadku ich
        poprawności przenosi użytkownika do odpowiedniego panelu (nauczyciela lub ucznia),
  \item zakładki \emph{Rejestracja}, umożliwiającej utworzenie nowego konta po podaniu imienia,
        nazwiska, loginu i hasła oraz wybraniu roli (ucznia lub nauczyciela),
  \item komunikatów zwrotnych, informujących o błędnych danych logowania lub o problemach
        z walidacją formularza rejestracji.
\end{itemize}

\figscreen{ekran_logowanie.png}{0.75}{Ekran logowania i rejestracji.}{rys:logowanie}

% =====================================================================
\section{Pulpit nauczyciela}
\label{sec:pulpit}

Po zalogowaniu na konto nauczyciela wyświetlany jest pulpit (rysunek~\ref{rys:pulpit}),
stanowiący przegląd najważniejszych informacji. Zawiera on:

\begin{itemize}
  \item kafelki podsumowujące liczbę testów, uczniów i grup oraz średni wynik,
  \item listę testów nauczyciela z szybkim przejściem do ich statystyk,
  \item panel ostatnich wyników uczniów,
  \item boczny panel nawigacji umożliwiający przejście do pozostałych modułów: testów, oceniania,
        grup i statystyk.
\end{itemize}

\figscreen{ekran_pulpit.png}{0.95}{Pulpit nauczyciela.}{rys:pulpit}

% =====================================================================
\section{Lista testów i edytor testu}
\label{sec:edytor}

Widok \emph{Testy} prezentuje wszystkie testy nauczyciela w formie tabeli z informacjami o nazwie
i statusie widoczności (aktywny/nieaktywny), przedmiocie, liczbie pytań, czasie, przypisanych grupach
i liczbie podejść. Status można przełączać jednym kliknięciem, a przy każdym teście dostępne są
akcje: statystyki, edycja, duplikowanie oraz usunięcie. Z poziomu tego widoku nauczyciel tworzy nowy
test lub otwiera istniejący w edytorze (rysunek~\ref{rys:edytor}).

Edytor testu pozwala określić ustawienia testu (nazwę, przedmiot wpisywany dowolnie lub wybierany
z listy, limit czasu, próg zaliczenia, liczbę podejść, widoczność dla uczniów oraz przypisane grupy)
i zbudować listę pytań. Dla każdego pytania można wybrać typ (jednokrotny wybór, wielokrotny wybór
lub pytanie otwarte), wpisać treść, ustalić liczbę punktów oraz - w pytaniach zamkniętych -
warianty odpowiedzi i poprawne odpowiedzi. Edytor waliduje kompletność testu przed zapisem (m.in.
obecność poprawnej odpowiedzi i przypisanej grupy). Jeżeli test ma już rozwiązania, edycja pytań
jest blokowana, o czym informuje stosowny komunikat.

\figscreen{ekran_edytor.png}{0.95}{Lista testów nauczyciela ze statusem widoczności i akcjami.}{rys:listaTestow}

\figscreen{ekran_testy.png}{0.95}{Edytor testu wraz z ustawieniami i listą pytań.}{rys:edytor}

% =====================================================================
\section{Ocenianie pytań otwartych}
\label{sec:ocenianie}

Moduł \emph{Ocenianie} (rysunek~\ref{rys:ocenianie}) gromadzi podejścia uczniów zawierające pytania
otwarte oczekujące na ręczną ocenę. Dla każdego zgłoszenia nauczyciel widzi:

\begin{itemize}
  \item dane ucznia oraz nazwę testu i datę oddania,
  \item treść pytania otwartego oraz odpowiedź wpisaną przez ucznia,
  \item podpowiedź w postaci odpowiedzi wzorcowej,
  \item pole do wpisania przyznanych punktów w dozwolonym zakresie (z walidacją zakresu).
\end{itemize}

Po przyznaniu punktów wszystkim pytaniom otwartym i zapisaniu oceny wynik ucznia jest automatycznie
przeliczany, a zgłoszenie znika z kolejki.

\figscreen{ekran_ocenianie.png}{0.9}{Ocenianie pytań otwartych.}{rys:ocenianie}

% =====================================================================
\section{Zarządzanie grupami}
\label{sec:grupy}

Widok \emph{Grupy} (rysunek~\ref{rys:grupy}) prezentuje grupy uczniów w formie kart. Każda karta
zawiera nazwę i opis grupy, listę przypisanych uczniów oraz listę testów udostępnionych grupie.
Przycisk zarządzania otwiera okno, w którym nauczyciel zaznacza uczniów należących do grupy.
Dodatkowo, z poziomu tego widoku można utworzyć nową grupę.

\figscreen{ekran_grupy.png}{0.95}{Widok zarządzania grupami uczniów.}{rys:grupy}

% =====================================================================
\section{Statystyki testu i eksport wyników}
\label{sec:statystyki}

Moduł \emph{Statystyki} (rysunek~\ref{rys:statystyki}) prezentuje analizę wyników wybranego testu.
Zawiera:

\begin{itemize}
  \item wskaźnik zdawalności oraz liczbę uczniów, łączną liczbę podejść i średni wynik,
  \item informację o najtrudniejszym pytaniu,
  \item wykres poprawności w rozbiciu na poszczególne pytania,
  \item tabelę z wszystkimi podejściami uczniów (najlepsze podejście oznaczone gwiazdką) wraz ze
        statusem (zaliczony, niezaliczony, oczekujący na ocenę, próba oszustwa) oraz możliwością
        ręcznej edycji oceny każdego podejścia,
  \item przycisk eksportu wyników do pliku CSV.
\end{itemize}

Wskaźniki zbiorcze (zdawalność, średni wynik) liczone są z najlepszego podejścia każdego ucznia.

\figscreen{ekran_statystyki.png}{0.95}{Statystyki wyników testu.}{rys:statystyki}

% =====================================================================
\section{Lista testów ucznia}
\label{sec:testyUcznia}

Po zalogowaniu na konto ucznia wyświetlana jest lista testów przypisanych do jego grup
(rysunek~\ref{rys:testyUcznia}), z podziałem na testy do zrobienia oraz ukończone. Każdy test
prezentowany jest w formie karty z informacją o przedmiocie, liczbie pytań, czasie, liczbie
wykorzystanych podejść oraz - dla testów rozwiązanych - z osiągniętym wynikiem i możliwością jego
podglądu lub ponownego podejścia (jeśli limit nie został wyczerpany, a test jest nadal aktywny).
Uczeń widzi wyłącznie testy aktywne oraz te nieaktywne, które już rozwiązał (zachowując dostęp do
swojego wyniku); po dezaktywacji testu niewykorzystane podejścia nie są już dostępne.

\figscreen{ekran_testy_ucznia.png}{0.95}{Lista testów przypisanych uczniowi.}{rys:testyUcznia}

% =====================================================================
\section{Rozwiązywanie testu}
\label{sec:rozwiazywanie}

Ekran rozwiązywania testu (rysunek~\ref{rys:rozwiazywanie}) prezentuje pojedyncze pytanie wraz
z wariantami odpowiedzi lub polem na odpowiedź otwartą. W górnej części znajduje się licznik czasu,
a u dołu - nawigacja umożliwiająca przechodzenie między pytaniami oraz pasek postępu. Po upływie
czasu test jest automatycznie kończony. Przy oddawaniu testu z pominiętymi pytaniami wyświetlane
jest ostrzeżenie. Opuszczenie okna testu jest traktowane jako próba oszustwa.

\figscreen{ekran_rozwiazywanie.png}{0.95}{Ekran rozwiązywania testu.}{rys:rozwiazywanie}

% =====================================================================
\section{Wynik testu}
\label{sec:wynik}

Po zakończeniu testu uczeń widzi ekran wyniku (rysunek~\ref{rys:wynik}), prezentujący osiągnięty
procent, liczbę zdobytych punktów oraz status (zaliczony, niezaliczony lub oczekujący na ocenę
pytań otwartych). Poniżej znajduje się szczegółowy przegląd odpowiedzi - dla każdego pytania
wskazane są odpowiedzi poprawne oraz wybrane przez ucznia.

\figscreen{wynik_testu.png}{0.95}{Ekran wyniku testu wraz z przeglądem odpowiedzi.}{rys:wynik}
      % 8
\chapter{Harmonogram realizacji projektu}
\label{cha:harmonogram}

W tym rozdziale przedstawiono harmonogram realizacji projektu, obrazujący poszczególne etapy pracy
oraz ich ramy czasowe. Do jego wizualizacji wykorzystano diagram Gantta
(rysunek~\ref{rys:gantt}), który pozwala w czytelny sposób przedstawić planowanie i nakładanie się
poszczególnych faz prac.

\begin{itemize}
  \item \textbf{Planowanie i projektowanie} -faza początkowa obejmująca analizę wymagań oraz
        zaprojektowanie ogólnej struktury aplikacji i schematu bazy danych.
  \item \textbf{Konfiguracja środowiska i bazy danych} -przygotowanie narzędzi, utworzenie
        projektu oraz konfiguracja połączenia z bazą danych i pierwszych migracji.
  \item \textbf{Implementacja modułów podstawowych} -rozwój podstawowych funkcji: logowania
        i rejestracji, zarządzania grupami oraz tworzenia testów.
  \item \textbf{Implementacja modułów zaawansowanych} -ocenianie odpowiedzi, obsługa wielu
        podejść, wykrywanie oszustw oraz statystyki i eksport wyników.
  \item \textbf{Testy i poprawki} -weryfikacja poprawności działania aplikacji oraz eliminacja
        napotkanych błędów.
  \item \textbf{Tworzenie dokumentacji} -przygotowanie dokumentacji, prowadzone częściowo
        równolegle z pracami implementacyjnymi.
\end{itemize}

\begin{figure}[H]
  \centering
  \resizebox{\textwidth}{!}{%
  \begin{tikzpicture}[x=0.95cm, y=0.95cm, font=\footnotesize]
    % --- siatka tygodni ---
    \foreach \x in {0,1,2,3,4,5,6,7,8,9,10,11,12} {
      \draw[gray!25] (\x,0.3) -- (\x,-6.3);
    }
    % --- etykiety miesięcy ---
    \node[anchor=south] at (2,0.4)  {Kwiecień};
    \node[anchor=south] at (6,0.4)  {Maj};
    \node[anchor=south] at (10,0.4) {Czerwiec};
    % --- zadania: etykieta po lewej + pasek ---
    \node[anchor=east] at (-0.2,-0.3) {Planowanie i projektowanie};
    \fill[blue!55]  (0,-0.05) rectangle (2,-0.55);
    \node[anchor=east] at (-0.2,-1.3) {Konfiguracja środowiska i bazy};
    \fill[blue!55]  (1.5,-1.05) rectangle (3,-1.55);
    \node[anchor=east] at (-0.2,-2.3) {Moduły podstawowe};
    \fill[teal!60]  (3,-2.05) rectangle (7,-2.55);
    \node[anchor=east] at (-0.2,-3.3) {Moduły zaawansowane};
    \fill[teal!60]  (6,-3.05) rectangle (10,-3.55);
    \node[anchor=east] at (-0.2,-4.3) {Testy i poprawki};
    \fill[orange!70] (9,-4.05) rectangle (12,-4.55);
    \node[anchor=east] at (-0.2,-5.3) {Dokumentacja};
    \fill[purple!55] (8,-5.05) rectangle (12,-5.55);
  \end{tikzpicture}%
  }
  \caption{Diagram Gantta przedstawiający harmonogram realizacji projektu.\label{rys:gantt}}
\end{figure}

\section{Napotkane trudności w trakcie realizacji projektu}
\label{sec:trudnosci}

Podczas realizacji projektu napotkano kilka trudności, które wymagały dodatkowego czasu i uwagi:

\begin{itemize}
  \item \textbf{Mapowanie dziedziczenia na bazę danych} -odwzorowanie trzech hierarchii klas
        (użytkownicy, pytania, odpowiedzi) w bazie wymagało prawidłowej konfiguracji strategii
        Table-per-Hierarchy oraz dyskryminatorów w Entity Framework Core.

  \item \textbf{Reguły usuwania powiązanych danych} -konieczne było staranne dobranie reguł
        kaskadowego i ograniczającego usuwania, tak aby z jednej strony usunięcie testu czyściło
        powiązane dane, a z drugiej - edycja testu z istniejącymi rozwiązaniami nie naruszała
        oddanych już odpowiedzi uczniów.

  \item \textbf{Obsługa wielu podejść i wyboru najlepszego wyniku} -wymagała rozróżnienia między
        zapisem nowego podejścia a aktualizacją oceny istniejącego oraz spójnego liczenia wyników
        po stronie ucznia i nauczyciela.

  \item \textbf{Wykrywanie prób oszustwa} -konieczne było prawidłowe rozpoznawanie utraty fokusu
        okna oraz odróżnienie własnych okien dialogowych aplikacji od faktycznego przełączenia się
        użytkownika na inny program.
\end{itemize}

\section{Repozytorium i system kontroli wersji}
\label{sec:repozytorium}

Kod źródłowy projektu jest przechowywany i zarządzany za pomocą systemu kontroli wersji Git. Dostęp
do niego możliwy jest poprzez publiczne repozytorium w serwisie GitHub. Korzystanie z systemu
kontroli wersji pozwoliło na śledzenie wszystkich zmian w kodzie oraz zarządzanie kolejnymi wersjami
projektu.

\noindent Adres repozytorium projektu:\\
\url{https://github.com/MichalNocek/ProjektPO2}
      % 9
\chapter{Instrukcja uruchomienia projektu}
\label{cha:uruchomienie}

\section{Wymagania systemowe}
\label{sec:wymaganiaUruch}

\begin{itemize}
  \item \textbf{System operacyjny:} Windows 10 lub Windows 11.
  \item \textbf{Środowisko uruchomieniowe:} platforma .NET (wersja zgodna z projektem) ze wsparciem WPF.
  \item \textbf{Serwer bazy danych:} PostgreSQL (zalecana wersja 13 lub nowsza).
  \item \textbf{Dodatkowe narzędzia:} środowisko programistyczne wspierające .NET (np. Visual Studio)
        oraz narzędzie do zarządzania bazą danych (np. pgAdmin).
\end{itemize}

\section{Konfiguracja bazy danych}
\label{sec:konfBaza}

Aplikacja łączy się z bazą danych PostgreSQL przy użyciu parametrów zdefiniowanych w klasie
\texttt{AppDbContext}: serwer \texttt{localhost}, port \texttt{5432}, baza \texttt{testownik},
użytkownik \texttt{postgres}. Przed pierwszym uruchomieniem należy zapewnić działający serwer
PostgreSQL z dostępem do tych parametrów (w razie potrzeby parametry połączenia można dostosować
w kodzie klasy \texttt{AppDbContext}).

Ręczne tworzenie bazy ani tabel nie jest wymagane - schemat zakładany jest automatycznie. Przy
starcie aplikacja wywołuje metodę \texttt{Database.Migrate()}, która tworzy bazę i nakłada wszystkie
migracje Entity Framework Core, a następnie klasa \texttt{DbSeeder} wypełnia ją danymi początkowymi
(przykładowy nauczyciel, uczniowie, grupy oraz testy).

\section{Konfiguracja środowiska i uruchomienie aplikacji}
\label{sec:uruchAplikacji}

Kolejne kroki uruchomienia:

\begin{enumerate}
  \item uruchomić serwer PostgreSQL i upewnić się, że jest dostępny na \texttt{localhost:5432};
  \item otworzyć projekt w środowisku Visual Studio (plik rozwiązania \texttt{ProjektPO2.slnx})
        i uruchomić aplikację (lub wykonać polecenie \texttt{dotnet run} w katalogu projektu);
  \item przy pierwszym uruchomieniu baza zostanie utworzona i zasilona danymi demonstracyjnymi;
  \item zalogować się przy użyciu jednego z kont demonstracyjnych.
\end{enumerate}

\section{Dane logowania testowego}
\label{sec:daneLogowania}

Po zasileniu bazy danymi początkowymi dostępne są konta demonstracyjne (hasło dla wszystkich:
\texttt{demo1234}):

\begin{itemize}
  \item nauczyciel - login \texttt{m.dabrowski};
  \item uczniowie - loginy \texttt{anna.k}, \texttt{piotr.n}, \texttt{maria.w} oraz kolejni.
\end{itemize}
     % 10
\chapter{Podsumowanie}
\label{cha:podsumowanie}

Celem projektu było stworzenie aplikacji wspierającej nauczyciela i ucznia w pełnym cyklu pracy
z testami - od ich ułożenia, przez udostępnienie i rozwiązanie, aż po ocenę i analizę wyników. Cel
ten został zrealizowany: powstała spójna, w pełni funkcjonalna aplikacja desktopowa napisana
w języku C\# na platformie .NET, z interfejsem graficznym opartym na technologii WPF i z relacyjną
bazą danych PostgreSQL obsługiwaną przez Entity Framework Core.

\section*{Zrealizowane funkcjonalności}

W ramach projektu udało się zrealizować wszystkie założone funkcje:

\begin{itemize}
  \item \textbf{Konta i role} - rejestracja (z wyborem roli) i logowanie użytkowników
        z rozróżnieniem ról nauczyciela i ucznia oraz odrębnymi panelami dla każdej z nich; każdy
        nauczyciel zarządza wyłącznie własnymi testami.
  \item \textbf{Grupy} - tworzenie grup i zarządzanie ich składem, z możliwością przynależności
        ucznia do wielu grup.
  \item \textbf{Testy} - tworzenie, edycja, duplikowanie i usuwanie testów złożonych z pytań
        jednokrotnego wyboru, wielokrotnego wyboru oraz pytań otwartych, wraz z walidacją
        kompletności testu oraz sterowaniem widocznością (test aktywny/nieaktywny).
  \item \textbf{Rozwiązywanie testów} - rozwiązywanie testów w ograniczonym czasie, z obsługą
        wielu podejść i wyborem najlepszego wyniku.
  \item \textbf{Ocenianie} - automatyczne ocenianie pytań zamkniętych (z punktami częściowymi)
        oraz ręczne ocenianie i edycja oceny dowolnego podejścia (z możliwością nadpisania punktów
        każdego pytania) z automatycznym przeliczaniem wyniku.
  \item \textbf{Bezpieczeństwo i integralność} - haszowanie haseł, unikalność loginów, wykrywanie
        prób oszustwa oraz spójność danych zapewniona przez klucze obce i migracje.
  \item \textbf{Analiza wyników} - statystyki testu z prezentacją wszystkich podejść uczniów oraz
        eksport wyników do pliku CSV.
\end{itemize}

W trakcie tworzenia projektu zastosowano zasady programowania obiektowego: dziedziczenie
(odwzorowane w bazie danych strategią Table-per-Hierarchy), enkapsulację oraz podział kodu na
niezależne, logiczne warstwy. Dzięki temu struktura aplikacji jest czytelna i łatwa w utrzymaniu.

\section*{Możliwości dalszego rozwoju}

Choć aplikacja jest w pełni funkcjonalna, można ją w przyszłości rozbudować, na przykład o:

\begin{itemize}
  \item losowanie kolejności pytań i odpowiedzi w celu dodatkowego utrudnienia ściągania,
  \item podgląd historii własnych podejść po stronie ucznia (obecnie uczeń widzi swój najlepszy
        wynik, a pełną historię podejść ma nauczyciel w statystykach),
  \item terminy dostępności testów (przedział dat, w którym test można rozwiązać),
  \item rozbudowane raporty i wykresy postępów uczniów w czasie.
\end{itemize}

Dodanie tych funkcji zwiększyłoby użyteczność systemu, jednak w obecnej formie projekt w pełni
spełnia założone wymagania i stanowi solidną podstawę do dalszej rozbudowy.
     % 11

%------------------------------------------------------------------------------
% Bibliografia (przed spisami)
%------------------------------------------------------------------------------

\cleardoublepage
\begingroup
  \renewcommand{\emph}[1]{\textit{#1}} % wyłącz ulem na czas bibliografii
  \addcontentsline{toc}{chapter}{Bibliografia}
  \bibliographystyle{plain}
  \bibliography{bibliografia}
\endgroup
\renewcommand{\emph}[1]{\uline{#1}} % przywróć ulem (opcjonalnie)

%------------------------------------------------------------------------------
% Spisy
%------------------------------------------------------------------------------

\cleardoublepage
\addcontentsline{toc}{chapter}{Spis rysunków}
\listoffigures

\cleardoublepage
\addcontentsline{toc}{chapter}{Spis tabel}
\listoftables

\cleardoublepage
\addcontentsline{toc}{chapter}{Spis listingów}
\lstlistoflistings

\end{document}
